\section{Conclusion}
\label{sec:conclusion}

At full training data, every trained method -- EMOS pooled, EMOS local,
and DRN -- beats TimesFM-3's CRPS, and a one-parameter variance-inflation
rescaling of the raw ensemble does too, with a fixed-climatological
variant achieving this using no training data at all
(Section~\ref{sec:results-full}). This advantage appears at very small
training sizes: EMOS pooled's raw CRPS crossing is at
\BpCrpsPooledCases{} cases (\BpCrpsPooledDays{} calendar days), and that
crossing becomes statistically significant, under day-blocking (our
primary block definition), from \BpCrpsPooledSigK{} cases
(\BpCrpsPooledSigDays{} days) onward for both EMOS variants -- the
station-blocked \(t\)-test does not reach significance anywhere in the
tested \(k=1\)--9 range, while the station-blocked Wilcoxon test does
from \(k=\BpCrpsPooledStationWilcoxonFirstSigK{}\) onward
(Section~\ref{sec:results-significance}). For a newly commissioned
station, this is the practical answer to this paper's title question on
CRPS: roughly \BpCrpsPooledSigDays{} days of paired ensemble/observation
history, not months, are enough to reach a statistically defensible CRPS
advantage over a frozen, zero-shot foundation model, and does so at a
lower inference cost too, on the qualitative comparison in
Section~\ref{sec:discussion-fewdays} (a real, measured EMOS fit time
against an unmeasured but qualitatively larger GPU inference cost for
TimesFM-3, not a quantified comparison on both sides). Coverage crosses
far more slowly
than CRPS -- \CoverageCrpsRatioPooled{}x later for pooled EMOS,
\CoverageCrpsRatioLocal{}x later for local EMOS -- so this practical
answer holds for CRPS specifically, not for calibration.

On the calibration-sharpness axis, TimesFM-3 holds a real advantage over
a trivial rescaling, and this is settled, not conditional: even a
leakage-optimistic multiplier fit directly on test-set coverage -- an
upper bound no legitimately deployable trivial rescaling could exceed --
needs \InflLeakyWidth{} K, still wider than TimesFM-3's \WidthTsfmFull{}
K at matched coverage, by a margin of \InflLeakyWidthDiffPoint{} K (95\%
CI [\InflLeakyWidthDiffCiLow{}, \InflLeakyWidthDiffCiHigh{}],
day-blocked bootstrap, Section~\ref{sec:discussion-trivial}) that
excludes zero entirely. This advantage is not an
artifact of comparing a leakage-free multiplier against a
leakage-optimistic ceiling; it survives even the most favorable trivial
baseline this paper could construct.

TimesFM-3's short-lead CRPS advantage reverses within the tested
lead-time range -- first sign flip at \FirstSignFlipHours{}~h, durable
reversal by \DurableCrossoverHours{}~h
(Section~\ref{sec:results-leadtime}) -- and the mechanism behind that
short-lead advantage has two independent parts, both now tested
separately. The observation-history part does not support an
information-asymmetry explanation: a persistence-augmented EMOS at
0--24~h closes only \PersistenceGapClosedPct{}\% of the gap toward
TimesFM-3. This result is not treated as confirming an architectural
explanation instead, since the test itself has three limitations that
keep it from cleanly distinguishing the two hypotheses -- a lead-pooled
EMOS fit, a single linear persistence term against TimesFM-3's much
richer context window, and a metric-dependent result (EMOS is
significantly closer to nominal coverage in the same bucket where
TimesFM-3 wins on CRPS) (Section~\ref{sec:discussion-mechanism}). The
ensemble-spread/calibration part has also now been tested directly:
TimesFM-3's covariate channel does accept a second, simultaneous input,
and the model is genuinely sensitive to that channel's real
informational content, not merely its presence
(Section~\ref{sec:covariates}, Section~\ref{sec:limitations}) -- but
supplying ensemble spread this way does not cleanly close the
calibration deficit. CRPS worsens from \CrpsTsfmFull{} K to
\CrpsTsfmTwoCov{} K; coverage moves closer to nominal, from
\CoverageTsfmFull{} to \CoverageTsfmTwoCov{} -- a genuine calibration
shift, not merely a wider interval -- but interval width also increases
substantially, from \WidthTsfmFull{} K to \WidthTsfmTwoCov{} K, the
expected cost of that shift under this paper's own standard of judging
sharpness jointly with calibration, not independent evidence of a worse
model (Section~\ref{sec:discussion-covariate-check}). The net effect is
a shift along the calibration-sharpness trade-off together with a worse
CRPS, not a clean improvement or a clean failure. Neither half of the
mechanism therefore supports the information-asymmetry framing -- but
neither confirms architecture as the explanation either: the channel
accepted the second input and was demonstrably sensitive to it, and did
shift its behavior in response, yet that shift did not convert into a
net improvement within this study's zero-shot, frozen-weights setup. The
short-lead advantage's mechanism remains an open interpretive question,
with architecture one candidate explanation among others, not a
confirmed one.

Two questions that were open earlier in this project are now both
resolved. EMOS's training/test ensemble-size mismatch is bounded and too
small to be a material contributor to the calibration-sharpness
comparison (Section~\ref{sec:limitations}). And whether ensemble spread
could be supplied to TimesFM-3 as a second covariate, and whether doing
so would change its calibration, has been tested directly
(Section~\ref{sec:discussion-covariate-check}): it can be supplied, but
doing so does not improve calibration in a way that supports the
information-asymmetry framing (Section~\ref{sec:discussion-mechanism}).
Neither resolution changes this paper's two central findings: on CRPS, a
frozen, zero-shot time series foundation model is
not competitive with a few days of station-specific training data, for
either a trained statistical postprocessing method or a trivial
one-parameter rescaling of the raw ensemble (though the exact training
size at which this becomes statistically significant, rather than a raw
point-estimate crossing, depends on block definition,
Section~\ref{sec:results-significance}); and on calibration-sharpness,
TimesFM-3's advantage over even a leakage-optimistic trivial rescaling is
real. The evidence for both is for one time series foundation model,
TimesFM-3; whether it generalizes to other pretrained time series models
(e.g.\ Chronos-2, Moirai-2) is untested. Both findings are also specific
to the target variable studied here, \(2\,\mathrm{m}\) temperature, where
EMOS's Gaussian predictive family is close to ideally specified; for
targets requiring non-Gaussian or target-specific reformulations, a
trained method's small-sample advantage should not be assumed to appear
as early.
